% !PRINTED 566
\weberpage{566}
%<<CONT
\[ \sum^{s} \varepsilon^{(\lambda+\mu)\,\mathrm{ind}\,s(t+1)} r^{s(t+1)} = \sum^{s} \varepsilon^{(\lambda+\mu)\,\mathrm{ind}\,s} r^s = (\varepsilon^{\lambda+\mu}, r), \]
und die Formel (4) ergiebt
\weq{5}{\frac{(\varepsilon^\lambda, r)(\varepsilon^\mu, r)}{(\varepsilon^{\lambda+\mu}, r)} = \psi_{\lambda,\mu}(\varepsilon),}
wenn zur Abk\"urzung
\weq{6}{\sum_{1,\,n-2}^{t} \varepsilon^{\mu\,\mathrm{ind}\,t - (\lambda+\mu)\,\mathrm{ind}(t+1)} = \psi_{\lambda,\mu}(\varepsilon)}
gesetzt und $\lambda + \mu$ durch $n - 1 = m$ nicht theilbar angenommen
wird. Diese Functionen $\psi$ sind dann aus den $(n-1)^{\mathrm{ten}}$ Einheitswurzeln
 $\varepsilon$ mit ganzzahligen Co\"efficienten zusammengesetzt und
sind also bekannt, wenn die $(n-1)^{\mathrm{ten}}$ Einheitswurzeln als bekannt
 vorausgesetzt werden. Es sind Zahlen des K\"orpers $R(\varepsilon)$.

Wir nehmen jetzt $\lambda$ durch $\mu$ theilbar an, ersetzen also $\lambda$
durch $\mu\lambda$ und setzen
\weq{7}{\alpha = \varepsilon^\mu,}
so dass $\alpha$ auch eine nicht primitive $(n-1)^{\mathrm{te}}$ Einheitswurzel,
z. B. eine $e^{\mathrm{te}}$ Einheitswurzel sein kann (wenn, wie oben, $m = ef$
ist und $\mu$ durch $f$ theilbar, etwa $= f$ genommen wird), wobei
jedoch der Fall $\alpha = 1$, also $e = 1$ auszuschliessen ist.
Dann setzen wir
\weq{8}{\sum_{1,\,n-2}^{t} \alpha^{\mathrm{ind}\,t - (\lambda+1)\,\mathrm{ind}(t+1)} = \psi_\lambda(\alpha);}
darin ist $\psi_\lambda(\alpha)$ eine Zahl des K\"orpers $R(\alpha)$, die mit H\"ulfe einer
Indextabelle leicht berechnet werden kann; $\lambda$ kann dabei nach
dem Modul $e$ genommen werden. Die Formel (5) giebt f\"ur diese
Annahme
\weq{9}{(\alpha^\lambda, r)(\alpha, r) = (\alpha^{\lambda+1}, r)\psi_\lambda(\alpha),}
und diese Formel gilt, so lange $\mu(\lambda+1)$ nicht durch $n - 1$
theilbar ist, also, wenn $\mu$ zu $m$ theilerfremd, und daher $\alpha$ eine
primitive $m^{\mathrm{te}}$ Einheitswurzel ist, f\"ur $\lambda = 1, 2 \ldots n-3$, und
wenn $f$ der gr\"osste gemeinschaftliche Theiler von $\mu$ und $m$ und
$m = ef$ ist, f\"ur $\lambda = 1, 2 \ldots e-2$.
Die Formel (3) giebt jetzt, wenn $\alpha^\lambda$ nicht $= 1$ ist,
\weq{10}{(\alpha^\lambda, r)(\alpha^{-\lambda}, r) = (-1)^{\mu\lambda} n.}
Die Formeln (9), (10) gen\"ugen, um nicht nur die $r$ selbst,
sondern auch die s\"ammtlichen Perioden $\eta$ durch Radicale zu
bestimmen, deren Grade die in $m$ aufgehenden Primzahlen sind,
und unter denen die Functionen $\psi_\lambda$ als die bekannten Gr\"ossen
vorkommen.
